\chapter{Viết và xuất bản với AI tạo sinh}

\subsection{Tóm tắt}

Chương 3 xem xét cách AI tạo sinh (generative AI) đang làm thay đổi hoạt động viết và công bố học thuật. Chương khảo sát vai trò của AI trong toàn bộ vòng đời viết, từ vượt qua tình trạng bí ý tưởng và cấu trúc hóa lập luận, đến trau chuốt văn phong, định dạng trích dẫn và chuẩn bị hồ sơ nộp bài. Trong khi AI đẩy nhanh việc soạn thảo và chỉnh sửa thông qua động não ý tưởng, biên tập văn phong, thậm chí cả phản biện đồng cấp (peer review) mô phỏng, chương này nhấn mạnh rằng sự giám sát của con người vẫn giữ vai trò then chốt đối với chiều sâu trí tuệ, tính chính xác của dữ kiện và sự nghiêm cẩn học thuật. Chương cũng đề cập các rủi ro của việc lệ thuộc quá mức, bao gồm trích dẫn do AI bịa ra (hallucinated citations) và sự xói mòn tiếng nói học thuật, đồng thời đưa ra các chiến lược thực tiễn để tích hợp AI một cách có trách nhiệm vào quy trình viết mà không làm tổn hại tính độc đáo (originality) và sự liêm chính.

\textbf{Từ khóa:} Viết học thuật, AI tạo sinh, xuất bản học thuật, quản lý trích dẫn, soạn thảo có hỗ trợ của AI

\end{quote}

Viết với AI tạo sinh đại diện cho một bước chuyển mình mang tính biến đổi trong thực hành học thuật. Các học giả không còn phải nhìn chằm chằm vào màn hình trống, loay hoay tìm ý tưởng hay cách mở đầu một bài báo hoặc một đề xuất xin tài trợ. Các công cụ AI tạo sinh có thể giúp vượt qua tình trạng bí ý tưởng khi viết (writer's block) bằng cách cung cấp cấu trúc, gợi ý cách diễn đạt và mở ra những góc nhìn mới để khai thác. Tác động của chúng vượt xa việc tạo văn bản đơn thuần: chúng là một công cụ đa năng xuyên suốt toàn bộ quy trình học thuật. Từ động não về các câu hỏi nghiên cứu và xây dựng lập luận, đến soạn thảo đề xuất xin tài trợ, tổng hợp tổng quan tài liệu và chuẩn bị các bài thuyết trình nghiên cứu, AI hỗ trợ ở mọi giai đoạn (Kulkarni et al., 2024). AI có thể giúp soạn tóm tắt hội nghị, phát triển phương pháp nghiên cứu, hình thành giả thuyết, tinh chỉnh khung lý thuyết và xác định các khoảng trống nghiên cứu tiềm năng trong một lĩnh vực (Wagner et al., 2022). Ngoài ra, các công cụ này hỗ trợ quản lý dự án học thuật bằng cách lập lộ trình nghiên cứu, tạo tài liệu giảng dạy và xây dựng thang đánh giá (rubric). Hơn cả một công cụ hỗ trợ viết, AI tạo sinh đang định hình lại cách giới học thuật suy nghĩ, lập kế hoạch và thực hiện công việc của mình, làm thay đổi căn bản quy trình nghiên cứu và xuất bản (Dwivedi et al., 2023).

Bất chấp những tiến bộ, AI tạo sinh vẫn chưa đạt đến mức chỉ với một prompt (câu lệnh) đã có thể tạo ra một bài báo học thuật hoàn chỉnh đạt chất lượng công bố. Đã có nhiều nỗ lực tự động hóa việc viết học thuật, nổi bật nhất có lẽ là mô hình STORM của Stanford, nhưng bất kỳ ai từng thử nghiệm các hệ thống như vậy đều nhanh chóng nhận ra những hạn chế của chúng (Kulkarni et al., 2024). Dù AI có thể tạo ra văn bản mạch lạc và thậm chí chèn trích dẫn, nghiên cứu và xuất bản học thuật đòi hỏi nhiều hơn rất nhiều so với việc ghép các từ ngữ theo một khuôn dạng có cấu trúc. Một bài báo được xây dựng tốt là sản phẩm của phân tích chặt chẽ, sự tương tác phản biện với tài liệu, độ chính xác phương pháp luận và sự sáng tạo trí tuệ, những yếu tố mà AI, ít nhất là hiện nay, không thể tự mình tái tạo (Milano et al., 2023). Dù có thể sẽ đến ngày AI đóng vai trò lớn hơn trong việc tạo ra nghiên cứu, các học giả vẫn cần tự thực hiện nghiên cứu, diễn giải kết quả và xây dựng lập luận với chiều sâu và tính độc đáo (Dwivedi et al., 2023).

Chương này khám phá các cách thức khác nhau mà AI tạo sinh có thể hỗ trợ việc viết học thuật, đồng thời nhấn mạnh sự cần thiết của giám sát từ con người. Chúng tôi bắt đầu với quá trình soạn thảo, xem xét cách AI có thể hỗ trợ động não, cấu trúc hóa và trau chuốt văn bản, đồng thời bảo đảm các học giả giữ quyền kiểm soát đối với đóng góp trí tuệ của mình. Tiếp theo, chúng tôi chuyển sang trích dẫn và tài liệu tham khảo, thảo luận cách AI có thể giúp định dạng danh mục tài liệu tham khảo, quản lý các kiểu trích dẫn và đối chiếu nguồn, đồng thời đề cập các rủi ro tiềm ẩn như tài liệu tham khảo do AI bịa ra (hallucination) (Milano et al., 2023). Cuối cùng, chúng tôi khám phá các chiến lược tối ưu hóa việc xuất bản, bao gồm cách AI có thể giúp điều chỉnh bản thảo phù hợp với hướng dẫn của tạp chí, phân tích các công bố gần đây để đáp ứng kỳ vọng của ban biên tập và hỗ trợ soạn các tài liệu nộp bài như tóm tắt và thư giới thiệu bài (cover letter). Xuyên suốt, chúng tôi nhấn mạnh tầm quan trọng của phán đoán con người trong việc bảo đảm AI vẫn là công cụ nâng cao năng lực chứ không thay thế tư duy phản biện và sự chặt chẽ học thuật.

\section{Quy trình soạn thảo}

Viết lách thường là một nhiệm vụ đầy thách thức, ngay cả với những học giả giàu kinh nghiệm. Hiện tượng bí ý tưởng (writer's block) nổi tiếng không chỉ xảy ra với sáng tác văn chương; nó cũng thường xuyên ảnh hưởng đến công việc học thuật. Bạn đã bao nhiêu lần nhìn chằm chằm vào màn hình trống, loay hoay không biết bắt đầu, viết tiếp hay kết thúc một bài báo nghiên cứu, một đề xuất tài trợ hay bất kỳ văn bản học thuật nào khác? Tin tốt là AI tạo sinh có thể giúp vượt qua những rào cản này (Grimes et al., 2023). Dù bạn đang soạn một bài báo, hình thành ý tưởng cho một đề xuất tài trợ nghiên cứu, hay xây dựng cấu trúc cho một luận án, các công cụ AI đều có thể hỗ trợ ở mọi giai đoạn của quá trình viết. Bạn không còn phải bị tê liệt trước viễn cảnh đáng ngại của một trang giấy trắng hay một lập luận dang dở, AI có thể cung cấp cấu trúc, gợi ý ý tưởng và trau chuốt văn bản của bạn một cách hiệu quả.

Tất nhiên, cách AI hỗ trợ việc soạn thảo sẽ khác nhau tùy theo loại văn bản học thuật. Viết một bài báo nghiên cứu khác với viết một đề xuất tài trợ, và cả hai có thể khác với việc soạn một chương sách hay một bản tóm tắt hội nghị. Tuy nhiên, các nguyên tắc cơ bản vẫn như nhau: AI có thể đóng vai trò là đối tác động não, công cụ tạo văn bản ban đầu và công cụ trau chuốt. Dù các kỹ thuật cụ thể có thể khác nhau, điểm cốt lõi cần ghi nhớ là AI tạo sinh không thay thế tư duy phản biện và chuyên môn của tác giả, mà là một công cụ giúp tinh giản và nâng cao quá trình viết (Cotton et al., 2023).

Trước hết cần làm rõ rằng, dù AI là một công cụ viết mạnh mẽ, việc dựa vào nó để tạo ra toàn bộ các phần từ đầu là điều không nên. Sự cám dỗ chỉ cần nhập một prompt (câu lệnh) như ``Hãy viết cho tôi một đề xuất nghiên cứu hoàn chỉnh về X'' và nhận lại một phản hồi có cấu trúc tốt, trông có vẻ mạch lạc, có thể rất lớn. Trong một số trường hợp, một prompt (câu lệnh) đủ chi tiết có thể cho ra văn bản tuân theo các hướng dẫn mong đợi. Tuy nhiên, dù kết quả như vậy có thể đọc được, nó khó có thể thực sự hữu ích. Các bản thảo do AI tạo ra thường thiếu chiều sâu, tính độc đáo và sự tương tác có ý nghĩa với chủ đề (Holmes, 2023). Thay vì tinh giản quá trình viết, cách tiếp cận này rốt cuộc có thể phản tác dụng, đòi hỏi phải chỉnh sửa, tái cấu trúc và trau chuốt nhiều để đạt chuẩn mực học thuật. Một chiến lược hiệu quả hơn là coi AI như một trợ lý giúp nâng cao và hoàn thiện các ý tưởng gốc, thay vì một người viết thuê thay thế nỗ lực trí tuệ.

Thay vì coi AI là sự thay thế cho việc viết, một cách tiếp cận hiệu quả hơn là bắt đầu bằng giai đoạn động não, như chúng ta đã thảo luận ở một chương trước. Tương tác với AI như một đối tác tư duy cho phép bạn tinh chỉnh ý tưởng, khám phá các góc nhìn thay thế và xây dựng một kế hoạch có cấu trúc trước khi soạn thảo. Hãy chọn công cụ phù hợp nhất với nhu cầu của bạn, dù là ChatGPT, Claude, hay một hệ thống AI tạo sinh khác, và cung cấp càng nhiều ngữ cảnh liên quan càng tốt. Hãy xem AI như một trợ lý nghiên cứu tài giỏi hoặc thậm chí một đồng nghiệp cộng tác, một đối tượng tuy đôi khi mắc lỗi (điều chúng ta sẽ đề cập sau), nhưng luôn sẵn sàng thảo luận, hào hứng trau chuốt công việc của bạn và có khả năng tạo ra những hiểu biết mới giúp làm sắc nét nghiên cứu của bạn (Stokel-Walker, 2023). Bằng cách chủ động tương tác với AI thay vì thụ động chấp nhận kết quả đầu ra của nó, bạn có thể tối đa hóa tiềm năng của nó như một công cụ nâng cao chất lượng viết học thuật.

Hãy cung cấp cho công cụ AI tạo sinh càng nhiều ngữ cảnh càng tốt. Bạn cung cấp càng nhiều thông tin, kết quả đầu ra càng chính xác và phù hợp. Hãy phác thảo điều bạn muốn viết, nêu rõ luận điểm của bạn và chỉ rõ các lập luận hoặc chủ đề chính. Nếu bạn đã có một bản thảo dang dở, hãy đưa vào, bằng cách đính kèm hoặc dán trực tiếp vào công cụ AI. Hãy đề cập mọi tài liệu, lý thuyết hay khung khái niệm liên quan mà bạn dự định đối thoại để AI có thể điều chỉnh các gợi ý cho phù hợp (Barros et al., 2023). AI hiệu quả nhất khi có đầu vào phong phú để làm việc, thay vì bị yêu cầu tạo nội dung từ con số không. Bằng cách coi nó như phần mở rộng của quá trình tư duy của chính bạn thay vì một lối tắt, bạn có thể bảo đảm rằng những đóng góp của nó phù hợp hơn với mục tiêu học thuật của bạn.

Tuy nhiên, trong quá trình động não này, điều quan trọng là phải nhận ra một hạn chế phổ biến của AI tạo sinh: nó có xu hướng đồng tình với bạn quá dễ dàng. Nếu bạn đưa ra một ý tưởng, AI thường sẽ đáp lại bằng \textbf{``Đó là một ý tưởng tuyệt vời!''} hoặc tạo ra những lý lẽ ủng hộ mà không thực sự phản biện lập luận của bạn. Điều này có thể tạo ra cảm giác được xác nhận sai lệch, khiến cách tiếp cận của bạn dường như không thể bị công kích. Tuy nhiên, việc viết học thuật đích thực phát triển nhờ sự chặt chẽ trí tuệ và thử thách. Để cân bằng xu hướng này, bạn phải chủ động phản bác lại những lời khẳng định của AI. Hãy yêu cầu rõ ràng nó phê bình các ý tưởng của bạn, chỉ ra các phản biện, hoặc nêu bật những điểm yếu tiềm tàng trong lập luận của bạn. Chẳng hạn, bạn có thể ra prompt (câu lệnh): \textbf{``Hãy thách thức luận điểm của tôi bằng cách chỉ ra những lỗ hổng trong lập luận của tôi hoặc những tài liệu hiện có mâu thuẫn với khẳng định của tôi.''} Thực hành này giúp bảo đảm rằng AI không chỉ củng cố quan điểm của bạn mà còn chủ động hỗ trợ việc hoàn thiện và tăng cường công trình của bạn.

Sau khi đã khám phá các ý tưởng của mình thông qua động não, bước tiếp theo là cấu trúc hóa bài viết. Trước khi đi vào các đoạn văn hoàn chỉnh, hãy cân nhắc khung tổng thể của bài báo, đề xuất hoặc chương sách. Bạn dự kiến độ dài là bao nhiêu? Bạn đang viết cho một tạp chí cụ thể, một đơn vị tài trợ hay một nhóm độc giả học thuật nào? Các phần chính và phần phụ là gì? Các công cụ AI có thể đặc biệt hữu ích ở giai đoạn này. Bạn có thể yêu cầu chúng đề xuất một dàn ý dựa trên thông tin bạn đã cung cấp: ``Given my thesis and main arguments, propose a structured outline for a journal article'' (Dựa trên luận điểm và các lập luận chính của tôi, hãy đề xuất một dàn ý có cấu trúc cho một bài báo khoa học) hoặc ``Based on my research topic, suggest a logical chapter breakdown.'' (Dựa trên chủ đề nghiên cứu của tôi, hãy gợi ý một cách phân chia chương hợp lý.) Dù không bao giờ nên làm theo các dàn ý do AI tạo ra một cách mù quáng, chúng vẫn có thể là điểm khởi đầu hữu ích, thúc đẩy bạn cân nhắc những cấu trúc tổ chức mà có thể bạn chưa nghĩ tới. Bước này giúp bảo đảm bản thảo của bạn có một lộ trình rõ ràng trước khi bạn bắt đầu quá trình viết thực sự.

Ở giai đoạn này, điều thiết yếu là tinh chỉnh các đầu ra do AI tạo ra càng nhiều càng tốt. Hãy phản biện các gợi ý của nó, chỉnh sửa khi cần thiết, và bảo đảm cấu trúc phù hợp với tầm nhìn của bạn. Đừng chấp nhận cho đến khi bạn hoàn toàn hài lòng với cách tổ chức các chương và tiểu mục (nếu viết một bài báo) hoặc cấu trúc tổng thể của tài liệu học thuật của mình. Dù bạn luôn có thể xem lại và điều chỉnh sau, việc thiết lập một khung rõ ràng ngay bây giờ sẽ tiết kiệm đáng kể thời gian và công sức về lâu dài. Tác phẩm của bạn càng có cấu trúc và mạch lạc ở giai đoạn này thì quá trình viết và chỉnh sửa càng suôn sẻ.

Khi đã hoàn tất cấu trúc, bạn có thể chuyển sang giai đoạn tiếp theo: viết. Một trong những chiến lược hiệu quả nhất vào lúc này là đơn giản bắt đầu, mà không suy nghĩ quá nhiều. Viết học thuật thường bắt đầu như một tập hợp các ý tưởng rời rạc, các lập luận chưa hoàn chỉnh, hoặc các điểm liên kết lỏng lẻo. Mấu chốt là đặt được chữ lên trang giấy. Bạn không cần hướng tới sự hoàn hảo ở bản thảo đầu tiên; thay vào đó, hãy tập trung ghi lại suy nghĩ của mình, ngay cả khi chúng có vẻ chưa đầy đủ hay chưa trau chuốt. Thay vì chờ đợi một câu hoàn hảo xuất hiện, hãy viết ra bất cứ điều gì chợt đến. Dù đó là một luận điểm sơ bộ, một dàn ý, hay chỉ là vài thuật ngữ chính, việc đặt chữ lên trang giấy là bước đầu tiên. Bạn có thể làm điều này trực tiếp trong một công cụ AI tạo sinh hoặc trong một trình soạn thảo văn bản thông thường. Khi gặp trở ngại, dù là khó diễn đạt một khái niệm, thiếu một ví dụ, hay cảm thấy không chắc về một đoạn chuyển ý, hãy dùng dấu ngoặc vuông như ``{[} {]}'' để để lại một chỗ giữ chỗ và chèn một prompt (câu lệnh) ngắn nhờ AI hỗ trợ. Chẳng hạn, bạn có thể viết: \textbf{{[}Complete this section with an example of X{]}} (Hoàn thành phần này bằng một ví dụ về X) hoặc \textbf{{[}Refine this paragraph to improve clarity{]}} (Tinh chỉnh đoạn này để tăng độ rõ ràng). Nếu bạn có một ý cụ thể trong đầu nhưng khó diễn đạt hiệu quả, bạn có thể cung cấp thêm hướng dẫn, chẳng hạn: \textbf{{[}Rewrite this section: What I mean to say is that my proposed theory challenges the conventional understanding of\ldots{]}} (Viết lại phần này: Điều tôi muốn nói là lý thuyết tôi đề xuất thách thức cách hiểu thông thường về\ldots).

Ở giai đoạn này, đừng lo về lỗi ngữ pháp hay những điểm chưa hoàn hảo về văn phong, AI sẽ xử lý chúng thay bạn. Một trong những lợi thế đáng kể nhất của AI tạo sinh là khả năng diễn giải và tinh chỉnh các bản thảo thô, ngay cả khi chúng chứa lỗi chính tả, câu chưa hoàn chỉnh hoặc cách diễn đạt vụng về. Miễn là bạn cung cấp đủ ngữ cảnh và nội dung, AI nhìn chung sẽ hiểu được ý định của bạn và tạo ra một phiên bản trau chuốt hơn cho văn bản của bạn. Điều này cho phép bạn tập trung hoàn toàn vào ý tưởng và lập luận mà không bị phân tâm bởi những chỉnh sửa bề mặt. Lần đầu tiên trong lịch sử học thuật, các học giả có thể ưu tiên phát triển nội dung trong khi giao phần kỹ thuật ngôn ngữ cho AI, giải phóng bạn khỏi gánh nặng tự biên tập liên tục trong quá trình soạn thảo.

Đáng chú ý, quá trình viết sẽ khác nhau ở mỗi người, một số người có thể thích tự soạn từng đoạn văn một cách độc lập trước khi dùng AI để tinh chỉnh, trong khi những người khác có thể viết theo cách rời rạc hơn, dùng AI để lấp các khoảng trống trong lúc viết. Cả hai cách tiếp cận đều hợp lệ, và mấu chốt là tìm ra một phương pháp phù hợp với quy trình làm việc của bạn. Một số người viết có thể tự soạn toàn bộ các phần rồi dùng AI để nâng cao độ rõ ràng, tính mạch lạc hoặc lập luận. Những người khác có thể viết các bản thảo thô xen kẽ với các phần hoàn thiện do AI tạo ra, dùng các prompt như ``\textbf{{[}Continue this argument with an example{]}}'' (Tiếp tục lập luận này bằng một ví dụ) hoặc ``\textbf{{[}Strengthen this paragraph by incorporating a counterpoint{]}}'' (Củng cố đoạn này bằng cách đưa vào một luận điểm phản biện). Không có một cách đúng duy nhất để dùng AI trong soạn thảo; đúng hơn, vấn đề là tìm ra sự cân bằng giúp bạn luôn tham gia tích cực, đồng thời bảo đảm AI đóng vai trò một cộng sự hữu ích chứ không phải sự thay thế cho nỗ lực trí tuệ của bạn.

Khi soạn thảo, điều thiết yếu là duy trì tính liên tục trong các tương tác với AI. Nhiều công cụ AI tạo sinh lưu giữ một mức độ nào đó bộ nhớ hội thoại, nghĩa là giữ quá trình viết của bạn trong cùng một phiên có thể cải thiện tính nhất quán. Nếu bạn thường xuyên bắt đầu các cuộc hội thoại mới, AI có thể coi mỗi prompt như một yêu cầu riêng lẻ, thiếu ngữ cảnh từ các lần trao đổi trước. Điều này đặc biệt quan trọng khi làm việc với một dự án dài hạn như bài báo khoa học, luận án hoặc đề xuất xin tài trợ. Một số công cụ AI đang dần cải thiện khả năng ghi nhớ chi tiết qua nhiều lần tương tác, nhưng hiện tại, tốt nhất là bảo đảm tính liên tục bằng cách ở lại trong cùng một luồng hội thoại bất cứ khi nào có thể. Tuy vậy, hãy lưu ý đến những hạn chế của AI, theo thời gian, ngay cả trong cùng một cuộc hội thoại, nó có thể bắt đầu quên các chi tiết trước đó. Khi cần thiết, hãy định kỳ nhắc lại các điểm chính để bảo đảm AI luôn đồng bộ với mục tiêu của bạn.

Bất kể bạn chọn phương pháp soạn thảo nào, điều cốt yếu là luôn \textbf{chủ động phản biện} đối với nội dung do AI tạo ra, đặc biệt ở những giai đoạn đầu. AI tạo sinh không thực sự ``hiểu'' chủ đề của bạn, mà chỉ dự đoán văn bản dựa trên các mẫu hình, nghĩa là nó có thể tạo ra những đầu ra chặt chẽ về cấu trúc nhưng yếu về mặt khái niệm. Nó có thể hiểu sai phong cách bạn muốn, đưa vào những lập luận không liên quan, hoặc tạo ra văn bản mang cảm giác chung chung, hời hợt. Nếu bạn không đánh giá nghiêm ngặt vài đoạn đầu tiên có sự hỗ trợ của AI, bạn có nguy cơ phải sửa chữa rất nhiều về sau. Thay vào đó, hãy dành thời gian trau chuốt cẩn thận các phần mở đầu, bảo đảm chúng phù hợp với lập luận, phong cách và đối tượng độc giả dự kiến của bạn. Bằng cách đặt ra một tiêu chuẩn cao ngay từ đầu, bạn có thể làm cho quá trình soạn thảo trôi chảy hơn và giảm thiểu nhu cầu chỉnh sửa lớn về sau.

Khi đã có bản thảo đầu tiên, bạn có thể bắt đầu tinh chỉnh cả nội dung lẫn văn phong. Ở giai đoạn này, AI đặc biệt hữu ích trong việc định hình giọng điệu bài viết. Dù bạn muốn văn bản phản ánh phong cách quen thuộc của mình, mô phỏng một giọng văn học thuật nhất định, hay cân bằng giữa tính dễ tiếp cận và tính chặt chẽ, AI đều có thể hỗ trợ trau chuốt. Một số công cụ AI cho phép bạn tải lên các văn bản mẫu để phân tích mẫu hình hành văn, trong khi những công cụ khác có thể điều chỉnh đầu ra dựa trên chỉ dẫn trực tiếp. Chẳng hạn, bạn có thể đưa ra prompt (câu lệnh): ``Hãy viết lại đoạn văn này theo giọng điệu học thuật trang trọng hơn'' hoặc ``Hãy làm cho đoạn này rõ ràng hơn nhưng vẫn giữ văn phong chuyên nghiệp.'' Tuy nhiên, AI không phải lúc nào cũng làm đúng ngay từ lần đầu. Việc tinh chỉnh văn phong đòi hỏi sự lặp lại: nếu một bản chỉnh sửa do AI tạo ra không đáp ứng kỳ vọng của bạn, hãy đưa ra phản hồi rõ ràng. Hãy cụ thể nhất có thể, ví dụ: ``Phiên bản này quá dài dòng; hãy cô đọng lại nhưng vẫn giữ nguyên lập luận chính.'' Hãy coi AI như một trợ lý có năng lực nhưng thiếu chính xác: cần được chỉ dẫn rõ ràng và điều chỉnh liên tục để đáp ứng đúng kỳ vọng của bạn.

Một tính năng mạnh mẽ khác của việc viết có sự hỗ trợ của AI là khả năng giúp tinh chỉnh và cấu trúc hóa lập luận một cách linh hoạt. Nếu cảm thấy một phần nào đó còn yếu hoặc thiếu chiều sâu, bạn có thể dùng AI để mở rộng phần đó bằng cách đưa ra prompt: ``Hãy củng cố lập luận này bằng cách bổ sung một góc nhìn từ {[}một lĩnh vực hoặc học giả cụ thể{]}'' hoặc ``Hãy xác định những phản bác tiềm tàng đối với nhận định này và đề xuất các luận điểm phản biện.'' Quy trình này có thể đặc biệt hữu ích khi xây dựng một khung lý thuyết phức tạp, vì AI có thể giúp làm nổi lên những ý tưởng mà ban đầu bạn chưa nghĩ tới. Tuy nhiên, điều quan trọng là phải giữ thái độ hoài nghi: các phần mở rộng do AI tạo ra đôi khi có thể đưa vào những điều thiếu chính xác, diễn giải sai các cuộc tranh luận học thuật, hoặc đưa ra những lập luận có vẻ hợp lý nhưng thiếu cơ sở vững chắc trong tài liệu học thuật (chúng ta sẽ quay lại vấn đề này). Hãy luôn đối chiếu các gợi ý do AI tạo ra với các nguồn thực tế để bảo đảm công trình của bạn vẫn chặt chẽ và có căn cứ vững vàng.

Đến giai đoạn này, bạn nên có một bản thảo khả dụng gồm các lập luận chính, các phần được cấu trúc rõ ràng và một văn phong đang dần hình thành. Tuy nhiên, trước khi chuyển sang các bước hiệu đính cuối cùng, cần cân nhắc cách AI có thể giúp ích về \textbf{tính trùng lặp, sự rõ ràng và độ chính xác}. Ngay cả những người viết giàu kinh nghiệm cũng có xu hướng vô tình lặp ý hoặc giải thích quá mức một số điểm. AI có thể hỗ trợ làm gọn văn bản bằng cách xác định những chỗ lập luận chồng chéo, câu văn phức tạp không cần thiết, hoặc phần giải thích có thể súc tích hơn. Bạn có thể đưa ra prompt: ``Hãy xác định các cụm từ thừa và đề xuất phiên bản súc tích hơn'' hoặc ``Hãy đơn giản hóa phần giải thích này nhưng vẫn giữ nguyên độ chính xác.'' Một số công cụ AI cũng có thể đánh giá cấu trúc câu, gợi ý cách cải thiện độ rõ ràng mà không hy sinh chiều sâu. Dù những chỉnh sửa này có giá trị, chúng luôn cần được xem xét một cách phản biện: AI không phải lúc nào cũng nhận ra các sắc thái của sự nhấn mạnh, các lựa chọn tu từ, hay sự lặp lại có chủ đích nhằm tạo hiệu quả. Hãy dùng nó như một công cụ để cải thiện, nhưng để phán đoán của chính bạn quyết định những thay đổi nào được giữ lại.

Vào thời điểm này, bạn cũng có thể tận dụng các công cụ AI để mô phỏng \textbf{quy trình phản biện đồng cấp (peer review)} trước khi nộp bài. Một kỹ thuật hữu ích là yêu cầu AI đóng vai nhiều người phản biện, mỗi người đánh giá bản thảo của bạn từ một góc nhìn khác nhau. Chẳng hạn, bạn có thể yêu cầu: ``Hãy đóng vai một người phản biện tạp chí và nhận xét phê bình bài báo này để đăng trên {[}tạp chí cụ thể{]}. Hãy chỉ ra các điểm yếu và đề xuất cải thiện.'' Nếu tạp chí có hướng dẫn nộp bài hoặc tiêu chí phản biện được công bố công khai, bạn thậm chí có thể đưa chúng vào prompt để nhận được nhận xét sát hơn với yêu cầu. Cách tiếp cận này có thể giúp phát hiện các lỗ hổng trong lập luận, làm nổi bật những chỗ cần được củng cố thêm, hoặc chỉ ra những điểm mà luận điểm của bạn có thể chưa rõ ràng với độc giả. Tuy nhiên, hãy lưu ý rằng các nhận xét do AI tạo ra, dù hữu ích, không thể thay thế phản hồi từ các chuyên gia thực thụ trong lĩnh vực của bạn. Đúng hơn, chúng nên được xem như một lớp tinh chỉnh bổ sung trước khi bạn xin ý kiến đồng nghiệp hoặc nộp bài cho quá trình phản biện đồng cấp chính thức.

Sau khi đã chỉnh sửa bản thảo dựa trên phản hồi do AI tạo ra, bạn có thể tiến thêm một bước bằng cách \textbf{lặp lại qua nhiều vòng rà soát}. Một sai lầm phổ biến là cho rằng một lần rà soát có sự hỗ trợ của AI là đủ để phát hiện mọi vấn đề. Để tránh điều này, hãy cân nhắc đưa bản thảo qua AI trong những bối cảnh khác nhau. Chẳng hạn, sau khi xử lý một loạt nhận xét, hãy mở một cuộc hội thoại AI mới và yêu cầu nó rà soát bản thảo đã chỉnh sửa như thể lần đầu nhìn thấy. Cách này giảm thiểu thiên lệch từ các phản hồi trước đó của AI và cho phép nhận được góp ý mới mẻ. Bạn cũng có thể dùng các mô hình AI khác nhau: điều mà công cụ này bỏ sót, công cụ khác có thể phát hiện. Nếu một vòng rà soát trước đã chỉ ra những điểm yếu về cấu trúc, vòng lặp tiếp theo có thể tập trung vào sự rõ ràng, tính súc tích hoặc sức mạnh lập luận. Cách tiếp cận theo lớp này giúp trau chuốt văn bản một cách tuần tự, bảo đảm rằng đến bản cuối cùng, công trình của bạn đạt mức hoàn thiện và chặt chẽ cao nhất có thể.

Ngoài việc phê bình và hoàn thiện bản thảo, AI còn có thể hỗ trợ \textbf{kiểm chứng thông tin và xác minh nguồn tài liệu}. Dù AI tạo sinh không thể thay thế một cuộc tổng quan tài liệu kỹ lưỡng, nó có thể giúp phát hiện các khoảng trống, điểm thiếu nhất quán, thậm chí cả lỗi trong các trích dẫn của bạn. Nếu bản thảo có những luận điểm cần được củng cố thêm, bạn có thể dùng prompt (câu lệnh) yêu cầu AI gợi ý các nguồn liên quan: ``Hãy xác định các bài báo học thuật quan trọng thảo luận về {[}chủ đề cụ thể{]}'' hoặc ``Hãy kiểm tra xem lập luận này có phù hợp với các nghiên cứu gần đây trong {[}lĩnh vực{]} hay không.'' Tuy nhiên, điều quan trọng là phải nhận thức rằng các mô hình AI, kể cả những mô hình được tích hợp chức năng tìm kiếm, có thể gây ra hiện tượng \emph{hallucinate} (ảo giác) về nguồn tài liệu, tức là bịa ra những trích dẫn nghe có vẻ hợp lý nhưng thực tế không tồn tại. Hãy luôn kiểm tra thủ công các tài liệu tham khảo bằng cách đối chiếu văn bản gốc và các cơ sở dữ liệu uy tín. Một thực hành tốt là sử dụng AI như công cụ để cải thiện chiến lược tìm kiếm của bạn, thay vì coi nó là nguồn trích dẫn trực tiếp. Cách làm này bảo đảm liêm chính học thuật và tránh việc dựa vào những tài liệu tham khảo không đáng tin cậy hoặc không tồn tại.

Khi bản thảo đã được trau chuốt và kiểm chứng thông tin, AI cũng có thể hỗ trợ \textbf{nâng cao tính dễ tiếp cận và mức độ thu hút}. Văn phong học thuật thường thiên về sự phức tạp, nhưng việc bảo đảm sự rõ ràng, đặc biệt đối với độc giả liên ngành, là điều then chốt. AI có thể giúp diễn đạt lại những đoạn dày đặc thông tin, đơn giản hóa phần giải thích, hoặc điều chỉnh mức độ chuyên môn phù hợp với đối tượng độc giả dự kiến. Chẳng hạn, bạn có thể ra prompt: ``Hãy viết lại đoạn văn này cho đối tượng độc giả rộng hơn mà vẫn giữ được tính chặt chẽ học thuật'' hoặc ``Hãy giải thích khái niệm này theo cách mà một người không chuyên trong lĩnh vực của tôi có thể hiểu được.'' Kỹ thuật này đặc biệt hữu ích khi chuẩn bị hồ sơ xin tài trợ, bản tóm lược nghiên cứu hướng đến công chúng, hoặc các công trình liên ngành mà sự rõ ràng là yếu tố quan trọng. Ngoài ra, nếu bạn cần điều chỉnh bài viết cho các định dạng khác nhau, chẳng hạn chuyển một bài báo khoa học thành bài thuyết trình hội nghị hoặc bài bình luận (op-ed), AI có thể hỗ trợ chuyển thể nội dung mà vẫn giữ nguyên các luận điểm cốt lõi.

Ở giai đoạn này, bạn cũng có thể cân nhắc sử dụng AI cho \textbf{hiệu đính lần cuối và điều chỉnh định dạng}. Dù các công cụ kiểm tra ngữ pháp và văn phong có hỗ trợ của AI không phải lúc nào cũng hoàn hảo, chúng có thể nhanh chóng phát hiện những lỗi chính tả bị bỏ sót, cách diễn đạt vụng về hoặc sự thiếu nhất quán về giọng văn. Các công cụ như ChatGPT, Grammarly hoặc Claude có thể chỉ ra những điểm yếu về cấu trúc, đề xuất cải thiện và bảo đảm tuân thủ các quy ước viết học thuật chính thức. Bạn có thể đưa ra những yêu cầu cụ thể cho AI, chẳng hạn ``Hãy kiểm tra văn bản này để tìm lỗi ngữ pháp và những điểm thiếu nhất quán'' hoặc ``Hãy bảo đảm bài viết này tuân theo hướng dẫn định dạng APA/Chicago/MLA.'' Điều này có thể tiết kiệm nhiều thời gian, đặc biệt khi chuẩn bị bản thảo nộp cho các tạp chí có yêu cầu định dạng nghiêm ngặt. Tuy nhiên, cũng như với mọi sản phẩm do AI tạo ra, sự giám sát thủ công là thiết yếu. Các chỉnh sửa tự động đôi khi có thể hiểu sai ngữ cảnh hoặc đưa vào những thay đổi ngoài ý muốn, vì vậy lần rà soát cuối cùng luôn nên do chính tác giả thực hiện để duy trì sự chính xác và tính xác thực.

Cuối cùng, AI có thể là một công cụ mạnh mẽ cho \textbf{tự đánh giá và cải thiện lặp lại}. Ngay cả sau khi đã hiệu đính và chỉnh định dạng, bạn vẫn có thể muốn đánh giá lại công trình của mình từ những góc nhìn khác nhau. Một kỹ thuật hiệu quả là dùng AI để mô phỏng nhiều kiểu độc giả khác nhau, chẳng hạn biên tập viên tạp chí, người thẩm định hồ sơ tài trợ, hoặc học giả thuộc các ngành lân cận. Bạn có thể ra prompt: ``Hãy đánh giá phần mở đầu này từ góc nhìn của người thẩm định hồ sơ tài trợ, xem xét độ rõ ràng và tầm quan trọng của nó'' hoặc ``Hãy phân tích lập luận này như thể bạn là một người phản biện (peer reviewer) đang tìm kiếm các điểm yếu logic.'' Bằng cách thực hiện nhiều vòng tự đánh giá có hỗ trợ của AI, bạn có thể chủ động giải quyết những vấn đề mà nếu không có thể bị chỉ ra trong quá trình phản biện chính thức. Cách tiếp cận lặp lại này bảo đảm bản nộp cuối cùng của bạn không chỉ được trau chuốt về ngôn ngữ và định dạng mà còn chặt chẽ về mặt trí tuệ, có cấu trúc tốt và được định vị chiến lược cho đối tượng độc giả mà bạn hướng đến.

Đây chỉ là những ví dụ về cách tương tác với các công cụ AI tạo sinh trong việc viết. Khi tích lũy thêm kinh nghiệm, bạn sẽ phát triển những phương pháp riêng để sử dụng các công cụ này như một trợ lý hiệu quả cho việc tư duy và viết. Với sự xuất hiện của các mô hình AI được thiết kế riêng cho việc viết học thuật, những khả năng này nhiều khả năng sẽ ngày càng tinh vi hơn, mang lại sự hỗ trợ tốt hơn trong việc xây dựng lập luận, trau chuốt văn phong và thích ứng với các quy ước của từng ngành.

Cũng đáng nói ở đây, dù chúng ta sẽ đi sâu hơn trong một chương khác, rằng khi bản thảo đã hoàn chỉnh, bạn nên kiểm tra đạo văn trước khi nộp. Nói chung, đạo văn không nên là mối lo ngại lớn khi AI tạo sinh được sử dụng đúng cách, vì các mô hình này tạo ra văn bản thay vì sao chép nguyên văn từ các nguồn hiện có. Tuy nhiên, rủi ro tăng lên khi sử dụng các công cụ AI truy xuất và tóm tắt nội dung bên ngoài. Nếu một mô hình AI lấy trực tiếp từ một bài báo trực tuyến hoặc cơ sở dữ liệu, kết quả đầu ra của nó có thể giống với tài liệu gốc, làm dấy lên những lo ngại về việc ghi nguồn không đúng cách.

Để bảo đảm tính nguyên bản, có thể sử dụng phần mềm phát hiện đạo văn. Một số công cụ phổ biến gồm trình kiểm tra đạo văn của Grammarly, Turnitin (thường được sử dụng trong các cơ sở học thuật) và iThenticate, được thiết kế riêng cho các nhà nghiên cứu và việc nộp bài cho tạp chí. Nhiều dịch vụ trong số này yêu cầu quyền truy cập thông qua tổ chức hoặc đăng ký trả phí, vì vậy bạn nên kiểm tra xem trường đại học hay tổ chức của mình có cung cấp quyền truy cập hay không trước khi mua giấy phép. Những biện pháp phòng ngừa này giúp duy trì liêm chính học thuật và bảo đảm rằng việc viết có hỗ trợ của AI vẫn phù hợp với các chuẩn mực đạo đức và nghề nghiệp.

Rõ ràng, việc sử dụng AI trong viết học thuật còn làm nảy sinh những cân nhắc đạo đức quan trọng khác. Mặc dù bản thân văn bản do AI tạo ra không cấu thành đạo văn, việc dựa vào AI để diễn đạt lại mà không trích dẫn đầy đủ có thể gây vấn đề về mặt đạo đức. Nếu AI hỗ trợ tóm tắt hoặc cấu trúc lại một lập luận bắt nguồn từ một nguồn hiện có, công trình gốc vẫn cần được ghi nhận. Xem AI như một trợ lý viết, chứ không phải sự thay thế cho các thực hành nghiên cứu và trích dẫn đúng đắn, là điều thiết yếu để duy trì liêm chính học thuật. Những mối quan ngại đạo đức này, cùng với các hệ quả rộng hơn của việc sử dụng AI trong môi trường học thuật, sẽ được khám phá chi tiết hơn ở phần sau của cuốn sách.

Tóm lại, quá trình soạn thảo với AI tạo sinh không nhằm thay thế nỗ lực trí tuệ của con người mà nhằm nâng cao hiệu quả, vượt qua chứng bí ý tưởng khi viết (writer's block) và trau chuốt văn bản học thuật. Bằng cách xem AI như một cộng tác viên tích cực chứ không phải một công cụ tạo nội dung thụ động, các học giả có thể tận dụng khả năng của nó để cấu trúc ý tưởng, phản biện lập luận và tinh gọn quy trình viết. Tuy nhiên, sự giám sát có tính phê phán vẫn là điều thiết yếu ở mọi giai đoạn, dù là khi kiểm chứng nội dung do AI tạo ra, tinh chỉnh các gợi ý về văn phong, hay bảo đảm tính nguyên bản và liêm chính học thuật được giữ vững. Khi đã có một bản thảo được cấu trúc tốt, các bước tiếp theo là quản lý trích dẫn và tài liệu tham khảo, trong đó AI có thể hỗ trợ định dạng danh mục tài liệu tham khảo và đối chiếu các nguồn, đồng thời tránh những cạm bẫy của các trích dẫn bị ``ảo giác'' (hallucination). Ngoài ra, AI có thể đóng vai trò trong việc tối ưu hóa chiến lược công bố bằng cách điều chỉnh bản thảo cho phù hợp với yêu cầu của tạp chí, soạn thảo các tài liệu nộp bài và điều chỉnh nghiên cứu cho những nhóm độc giả khác nhau. Khi chuyển sang các khía cạnh này của việc viết và công bố với AI, trọng tâm vẫn là tích hợp các công cụ này một cách có trách nhiệm, bảo đảm chúng đóng vai trò hỗ trợ chứ không thay thế chuyên môn học thuật.

\section{Trích dẫn và tài liệu tham khảo}

Việc trích dẫn nguồn đã trở nên dễ dàng hơn đáng kể nhờ các công cụ hỗ trợ bởi AI. Từ rất lâu trước khi AI tạo sinh xuất hiện, các phần mềm quản lý trích dẫn như Zotero, EndNote và Mendeley đã đơn giản hóa quá trình sắp xếp tài liệu tham khảo và định dạng trích dẫn đúng chuẩn. Những công cụ này ngày nay vẫn còn giá trị và đang dần tích hợp các khả năng của AI tạo sinh để nâng cao chức năng. Các mô hình AI tạo sinh như ChatGPT và Claude đã có thể hỗ trợ định dạng trích dẫn, quản lý danh mục tài liệu tham khảo, thậm chí đối chiếu chéo các tài liệu tham khảo. Với hướng dẫn phù hợp, AI có thể nhanh chóng định dạng trích dẫn theo APA, Chicago, MLA, Bluebook hoặc bất kỳ phong cách bắt buộc nào khác, giúp giảm đáng kể thời gian và công sức dành cho việc quản lý trích dẫn. Tuy nhiên, dù AI có thể hợp lý hóa các tác vụ trích dẫn, kết quả đầu ra của nó luôn cần được kiểm tra thủ công, vì đôi khi vẫn xảy ra sai sót, chẳng hạn như định dạng không đúng, thiếu thông tin hoặc thiếu nhất quán.

Tuy nhiên, về bản chất AI không đáng tin cậy khi tạo trích dẫn từ đầu. Một trong những vấn đề nghiêm trọng nhất là hiện tượng ảo giác của AI (AI hallucination) đã đề cập trước đó, tức là khi mô hình bịa đặt nguồn, tự tạo ra tên tác giả, tên tạp chí, năm xuất bản, thậm chí cả số trang. Nếu bạn yêu cầu AI cung cấp tài liệu tham khảo mà không đưa dữ liệu chính xác, nó có thể tạo ra những trích dẫn trông có vẻ hợp lệ nhưng hoàn toàn là hư cấu. Để tránh điều này, hãy luôn tìm và xác minh nguồn gốc trước khi dùng AI để định dạng trích dẫn. Thay vì yêu cầu AI trích dẫn các nguồn về một chủ đề nào đó, hãy cung cấp cho nó các tài liệu tham khảo thực tế. Nhiều công cụ AI hiện nay cho phép bạn tải trực tiếp tài liệu nguồn lên, giúp chúng trích xuất đúng các thông tin thư mục. Khi thông tin đã được xác minh, bạn có thể dùng AI để định dạng trích dẫn theo phong cách yêu cầu, dù là APA, Chicago, MLA, Bluebook hay phong cách khác. Ngoài ra, AI có thể giúp kiểm tra tính nhất quán của một danh sách trích dẫn, bảo đảm định dạng đồng nhất. Tuy nhiên, hãy luôn xem xét kỹ kết quả do AI tạo ra, vì nó vẫn có thể chứa những sai lệch nhỏ về định dạng. Nếu có lỗi, hãy sửa và hướng dẫn AI để các trích dẫn sau này phù hợp với phong cách bạn ưa thích.

AI cũng hữu ích trong việc quản lý toàn bộ danh mục tài liệu tham khảo. Nếu bạn có một danh sách tài liệu tham khảo, bạn có thể yêu cầu AI sắp xếp theo thứ tự bảng chữ cái, theo năm xuất bản, hoặc chuẩn hóa định dạng giữa các phong cách trích dẫn khác nhau. Điều này đặc biệt hữu ích khi hợp nhất các tài liệu tham khảo từ nhiều nguồn, bảo đảm tính nhất quán và tuân thủ các hướng dẫn phong cách cụ thể. Ngoài ra, AI có thể hỗ trợ phát hiện các thông tin còn thiếu, chẳng hạn tên tác giả chưa đầy đủ hoặc thiếu số trang, và đề xuất chỉnh sửa. Tuy nhiên, như với mọi nội dung do AI tạo ra, danh mục tài liệu tham khảo cuối cùng luôn cần được kiểm tra thủ công. AI đôi khi có thể bỏ sót thông tin quan trọng, sắp xếp sai thứ tự các mục, hoặc hiểu sai những khác biệt tinh tế giữa các phong cách trích dẫn. Dù các công cụ này có thể hợp lý hóa đáng kể quy trình, sự giám sát của con người vẫn là thiết yếu để bảo đảm tính chính xác và sự tuân thủ các chuẩn mực học thuật.

Đối với tác giả sách hoặc những người thực hiện các dự án nghiên cứu quy mô lớn cần có chỉ mục, AI cũng có thể là một công cụ hữu ích. Việc biên soạn chỉ mục thủ công là công việc tốn thời gian và đòi hỏi sự tỉ mỉ, nhưng AI có thể hỗ trợ tạo bản nháp sơ bộ. Bằng cách tải bản thảo lên, bạn có thể ra prompt (câu lệnh) để AI xác định các thuật ngữ then chốt, trích xuất những tên riêng và khái niệm quan trọng, hoặc sắp xếp các tham chiếu theo chủ đề hay danh mục cụ thể. Điều này có thể tạo ra một điểm khởi đầu hữu ích, tiết kiệm thời gian trong quá trình lập chỉ mục. Tuy nhiên, như với trích dẫn, việc lập chỉ mục đòi hỏi sự kiểm chứng cẩn thận. AI có thể gán sai số trang, bỏ sót các thuật ngữ quan trọng, hoặc không ưu tiên những khái niệm phù hợp nhất. Ngoài ra, nó có thể tạo ra một chỉ mục quá rộng hoặc cấu trúc chưa đủ chặt chẽ cho công trình học thuật chuyên sâu. Do đó, mọi chỉ mục do AI tạo ra nên được xem là bản nháp thô, cần được tinh chỉnh thủ công để bảo đảm tính chính xác, mạch lạc và phù hợp với văn bản.

Điểm mấu chốt ở đây là tìm được sự cân bằng phù hợp giữa việc dùng AI để tự động hóa các tác vụ tốn thời gian (và thường nhàm chán) với việc bảo đảm rằng nó không đưa vào những lỗi làm tổn hại đến công trình của bạn. AI có thể giải phóng đáng kể thời gian cho những khía cạnh đòi hỏi trí tuệ và sáng tạo hơn của nghiên cứu và viết lách, nhưng nó cần được hiệu chỉnh cẩn thận. Việc tìm ra sự cân bằng này thường là một quá trình thử và sai, tức là điều chỉnh cách bạn tích hợp AI vào quy trình làm việc cho đến khi nó phù hợp với nhu cầu cụ thể của bạn. Khi AI đã được tinh chỉnh theo sở thích của bạn, quy trình sẽ trở nên liền mạch hơn. Dù sự giám sát của con người vẫn là thiết yếu (như chúng tôi đã nhiều lần nhấn mạnh vì tầm quan trọng then chốt của nó), AI vẫn có thể đóng vai trò giá trị trong việc hợp lý hóa công việc học thuật của bạn.

\section{Tối ưu hóa việc xuất bản}

Nếu bạn đang chuẩn bị một bài báo, bước hợp lý tiếp theo sau khi viết xong là xuất bản. AI không nên thay thế các ưu tiên của cơ quan hoặc cá nhân bạn về nơi gửi công trình. Việc chọn đúng tạp chí hay nhà xuất bản phụ thuộc vào nhiều yếu tố, bao gồm nơi bạn muốn nghiên cứu của mình được biết đến, hệ số tác động của tạp chí, và mức độ phù hợp của công trình với độc giả của tạp chí đó. Mặc dù AI có thể hỗ trợ thu thập thông tin về các nơi xuất bản tiềm năng, bạn vẫn nên tham khảo ý kiến của đồng nghiệp, người hướng dẫn hoặc những người am hiểu bối cảnh xuất bản trong lĩnh vực của mình.

Khi đã xác định được tạp chí hoặc nhà xuất bản mục tiêu, AI có thể giúp tối ưu hóa chiến lược nộp bài. Nhiều tạp chí có các yêu cầu cụ thể về định dạng và văn phong, và AI có thể hỗ trợ bảo đảm bản thảo của bạn phù hợp với các hướng dẫn này. Bạn có thể tải lên các yêu cầu nộp bài của tạp chí và nhờ AI kiểm tra xem tài liệu của bạn có tuân thủ đúng hay không, đồng thời điều chỉnh các yếu tố như kiểu trích dẫn, số lượng từ và tiêu đề các phần cho phù hợp. Nếu bạn không chắc công trình của mình có nằm trong phạm vi của tạp chí hay không, bạn cũng có thể tải lên phần tóm tắt hoặc phần mở đầu của các số gần đây để phân tích các mô thức về văn phong, phương pháp luận hoặc cách đặt vấn đề. Điều này có thể mang lại những hiểu biết hữu ích về cách chỉnh sửa bản thảo để phù hợp hơn với kỳ vọng của tạp chí.

AI cũng có thể hỗ trợ soạn thảo các tài liệu nộp bài, chẳng hạn thư xin nộp bài (cover letter) hoặc phần tóm tắt. Nhiều tạp chí yêu cầu một thư xin nộp bài giải thích ý nghĩa của nghiên cứu và lý do nghiên cứu phù hợp với ấn phẩm của họ. Thay vì soạn từ đầu, bạn có thể cung cấp cho AI các thông tin chính, chẳng hạn câu hỏi nghiên cứu, những đóng góp chủ yếu và tên tạp chí, rồi nhờ nó tạo một bản nháp. Mặc dù các thư do AI tạo ra luôn cần được xem xét lại và cá nhân hóa, chúng có thể tiết kiệm thời gian bằng cách cung cấp một điểm khởi đầu có cấu trúc. Tương tự, AI có thể giúp trau chuốt phần tóm tắt để bảo đảm rõ ràng, súc tích và hấp dẫn, phù hợp với văn phong ưa chuộng của tạp chí.

Đối với những người nộp bài cho nhiều tạp chí, AI có thể hữu ích trong việc điều chỉnh một bài báo cho các nhóm độc giả khác nhau. Các tạp chí thường có phạm vi khác nhau, đòi hỏi tác giả phải điều chỉnh trọng tâm hoặc cách đặt vấn đề cho phù hợp với độc giả. AI có thể hỗ trợ quá trình này bằng cách gợi ý cách định vị lại các lập luận hoặc chỉnh sửa ngôn ngữ để phù hợp với những ưu tiên biên tập khác nhau. Tuy nhiên, như với mọi quy trình có AI hỗ trợ, sự giám sát của con người là thiết yếu: AI có thể đề xuất những thay đổi không hoàn toàn phù hợp với các lập luận cốt lõi của bài báo hoặc hiểu sai các quy ước của ngành.

Ngoài các bài báo tạp chí, AI cũng có thể hỗ trợ tối ưu hóa đề cương sách, bài gửi hội nghị và bản tóm lược nghiên cứu nhằm phổ biến rộng hơn. Dù là chuẩn bị một bài báo cho nền tảng truy cập mở, điều chỉnh bản thảo cho một ngành khác, hay tóm tắt các phát hiện cho công chúng, AI có thể giúp đơn giản hóa quá trình chuyển thể, đồng thời bảo đảm tính nhất quán của thông điệp.

\section{Kết luận}

AI tạo sinh (generative AI) đang định hình lại việc viết và xuất bản học thuật, nâng cao năng suất và tinh gọn quy trình làm việc. Từ việc động não và cấu trúc lập luận đến tinh chỉnh bản thảo, định dạng trích dẫn và tối ưu hóa hồ sơ nộp bài, AI đóng vai trò là một cộng tác viên có giá trị trong truyền thông học thuật. Tuy nhiên, AI không thay thế tư duy phản biện, tính độc đáo hay sự chặt chẽ về phương pháp luận, các học giả phải cân bằng giữa hiệu quả mà AI mang lại với sự giám sát của con người để duy trì liêm chính học thuật (academic integrity). Mặc dù AI có thể hỗ trợ quản lý trích dẫn (citation management), lập chỉ mục và tối ưu hóa bản thảo, các kết quả đầu ra của nó vẫn cần được xem xét cẩn trọng. Những rủi ro như trích dẫn bịa đặt (hallucinated citations), sự tiếp cận nông cạn với tài liệu và sự đồng nhất hóa diễn ngôn phải được chủ động giảm thiểu. Các tổ chức và nhà xuất bản ngày càng yêu cầu minh bạch trong việc sử dụng AI, khiến việc công bố thông tin một cách có trách nhiệm trở nên thiết yếu.

Nhìn về phía trước, vai trò của AI trong học thuật sẽ tiếp tục phát triển, làm nảy sinh những câu hỏi mới về quyền tác giả, sự ghi nhận nguồn và lao động trí tuệ. Thách thức không nằm ở việc có nên sử dụng AI hay không, mà là làm thế nào để tích hợp nó một cách có trách nhiệm. Bằng cách tận dụng các điểm mạnh của AI đồng thời bảo đảm sự chặt chẽ về trí tuệ, các học giả có thể nâng cao hiệu quả mà không làm tổn hại đến chiều sâu và tính độc đáo trong công trình của mình.
